\documentclass{beamer}

\usetheme{Madrid}
\usecolortheme{default}

% Packages for standard pdflatex compilation
\usepackage[T1]{fontenc}
\usepackage[utf8]{inputenc}
\usepackage{listings}
\usepackage{xcolor}
\usepackage{amsmath}

% Lean syntax highlighting setup for pdflatex
\lstdefinelanguage{lean}{
  keywords={theorem, def, variable, open, scoped, by, constructor, rintro, exact, have, rcases, calc, rw, apply, refine, intro, let, use, ext, dsimp, unfold, show, noncomputable, section, end},
  keywordstyle=\color{blue}\bfseries,
  ndkeywords={Type, Group, Subgroup, Set, Setoid, Quotient, Fintype},
  ndkeywordstyle=\color{teal}\bfseries,
  comment=[l]{--},
  commentstyle=\color{gray}\ttfamily,
  morecomment=[s]{/-}{-/},
  stringstyle=\color{red}\ttfamily,
  basicstyle=\ttfamily\scriptsize,
  breaklines=true,
  keepspaces=true,
  showstringspaces=false,
  % This block safely translates Lean's Unicode into LaTeX math equivalents
  literate=
    {∈}{{$\in$}}1
    {•}{{$\bullet$}}1
    {↔}{{$\leftrightarrow$}}1
    {∃}{{$\exists$}}1
    {∧}{{$\land$}}1
    {⁻}{{$^-$}}1
    {¹}{{$^1$}}1
    {⟨}{{$\langle$}}1
    {⟩}{{$\rangle$}}1
    {≃}{{$\simeq$}}1
    {→}{{$\to$}}1
    {⊆}{{$\subseteq$}}1
    {₁}{{$_1$}}1
    {₂}{{$_2$}}1
    {×}{{$\times$}}1
    {↦}{{$\mapsto$}}1,
}

\lstset{language=lean}

\title{Formalizing Lagrange's Theorem}
\subtitle{A Step-by-Step Approach in Lean 4}
\author{Formal Mathematics Presentation}
\date{\today}

\begin{document}

\begin{frame}
    \titlepage
\end{frame}

% --- SLIDE 1: The Statement ---
\begin{frame}[fragile]{The Statement We Want to Prove}
    \textbf{Lagrange's Theorem} is a fundamental result in group theory stating that for any finite group $G$ and subgroup $H$, the order of the subgroup divides the order of the group:
    
    \[ |G| = |G/H| \times |H| \]
    
    \vspace{0.5cm}
    In our Lean 4 formalization, this target statement looks like this:

\begin{lstlisting}
theorem lagrange_theorem 
  [Fintype G] [Fintype H] [Fintype (LeftCosets H)] :
  Fintype.card G = 
    Fintype.card (LeftCosets H) * Fintype.card H := by
    
    -- Our goal is to fill in this proof!
    ...
\end{lstlisting}
\end{frame}

% --- SLIDE 2: Proof Decomposition ---
\begin{frame}{Decomposing the Proof}
    To prove this elegantly in a proof assistant, we decompose the mathematics into three distinct milestones:
    
    \vspace{0.5cm}
    \begin{enumerate}
        \item \textbf{Cosets form a partition:} Show that for any two elements, their left cosets are either perfectly equal or completely disjoint. We do this by defining an equivalence relation on $G$.
        
        \vspace{0.3cm}
        \item \textbf{Equal cardinalities:} Prove that every left coset $gH$ has exactly the same cardinality as the subgroup $H$.
        
        \vspace{0.3cm}
        \item \textbf{$G$ as a product:} Show that the entire group $G$ can be written as the bijection (quotient space) $\times$ (subgroup).
    \end{enumerate}
\end{frame}

% --- SLIDE 3: Part 1 ---
\begin{frame}[fragile]{Part 1: Cosets are Equal or Disjoint}
    To show cosets partition $G$, we define a formal equivalence relation $a \sim b \iff a^{-1}b \in H$. The quotient space of this relation is exactly the set of disjoint left cosets.

\begin{lstlisting}
def leftRel (H : Subgroup G) : Setoid G where
  r a b := a⁻¹ * b ∈ H
  iseqv := {
    refl := by ...   -- Proves a⁻¹ * a = 1 ∈ H
    symm := by ...   -- Proves a⁻¹ * b ∈ H → b⁻¹ * a ∈ H
    trans := by ...  -- Proves a⁻¹ * b ∈ H ∧ b⁻¹ * c ∈ H → a⁻¹ * c ∈ H
  }

-- The set of all left cosets is the Quotient of this relation
def LeftCosets (H : Subgroup G) := Quotient (leftRel H)
\end{lstlisting}
    
    \vspace{0.2cm}
    \textit{This handles the first requirement: $G$ is now cleanly partitioned into $G/H$.}
\end{frame}

% --- SLIDE 4: Part 2 ---
\begin{frame}[fragile]{Part 2: Cosets Have the Same Cardinality}
    Next, we prove $|gH| = |H|$ by constructing an explicit, invertible mapping (a bijection, denoted by \texttt{≃}) between the subgroup $H$ and any arbitrary local coset $gH$.

\begin{lstlisting}
def LocalCoset (g : G) := {x : G // x ∈ g • (H : Set G)}

-- The equivalence proves they have the exact same number of elements
def local_coset_equiv (g : G) : H ≃ LocalCoset H g where
  -- Forward: Multiply an element h ∈ H by g
  toFun := by ... 
  
  -- Backward: Multiply an element x ∈ gH by g⁻¹
  invFun := by ...
  
  -- Proofs that these operations cancel each other out
  left_inv := by ...  -- g⁻¹ * (g * h) = h
  right_inv := by ... -- g * (g⁻¹ * x) = x
\end{lstlisting}
\end{frame}

% --- SLIDE 5: Part 3 ---
\begin{frame}[fragile]{Part 3: $G$ as a Quotient $\times$ Subgroup}
    We combine parts 1 and 2 to build a global bijection. Every element $g \in G$ is uniquely determined by which coset it belongs to, and its "position" inside that coset.

\begin{lstlisting}
noncomputable section

def groupEquivCosetProd (H : Subgroup G) : 
  G ≃ LeftCosets H × H where
  
  -- Take g ∈ G. Return its coset, and the "remainder" in H
  toFun g := ... 
  
  -- Take a coset q and h ∈ H. Return an element in G
  invFun p := ...
  
  -- Proofs of exact invertibility
  left_inv := by ...
  right_inv := by ...
\end{lstlisting}
\end{frame}

% --- SLIDE 6: Deducing the Formula ---
\begin{frame}[fragile]{Deducing the Formula}
    Because we did the heavy lifting of building the global bijection \texttt{groupEquivCosetProd}, the final theorem is deduced in just two steps using Lean's \texttt{Fintype} library.

\begin{lstlisting}
theorem lagrange_theorem 
  [Fintype G] [Fintype H] [Fintype (LeftCosets H)] :
  Fintype.card G = 
    Fintype.card (LeftCosets H) * Fintype.card H := by

  -- 1. Since G and (LeftCosets H × H) are bijective, 
  --    they have the same cardinality.
  rw [Fintype.card_congr (groupEquivCosetProd H)]

  -- 2. The size of a Cartesian product (A × B) is Size(A) * Size(B).
  simp only [Fintype.card_prod]
  
end
\end{lstlisting}
\end{frame}

\begin{frame}
    \begin{center}
        \Huge \textbf{Thank You!}
        
        \vspace{1cm}
        \large \textit{Questions?}
    \end{center}
\end{frame}

\end{document}